\section{The finite induction from fifteen to thirty-nine}
\label{ext:finite1539}
The same argument is repeated at each index, with different rational
bounds. Suppose that $n$ is the first index in this range for which
$C_n>\CE$. By Corollary~\ref{ext:ratio-attainment}, choose an unrestricted
positive closed ratio maximizer $P$ at this index. Its score is $C_n$,
and the same marginal satisfies $C_{n-1}(P)\le\CE<C_n$.
Consequently its contact predecessor allowance is zero. The positive
allowances retained in some calculations below are weaker valid inputs.
The law is never replaced when passing from a support calculation to a
moment calculation or to a conditional count probability.

There are two ways to exclude this law. For sufficiently large $\rho$,
a Fourier estimate directly bounds its score below $\CE$. In the
remaining moment intervals, support and Fourier estimates first give
necessary constraints. A finite moment calculation then gives a central
support gap and, on each probability panel, a joint bound on inward
displacement and conditional variance. The count calculation excludes
every distribution obeying those constraints. Its domain is larger than
the set of distributions, which is harmless in this direction.

\subsection{Moment intervals and the order of their inputs}
Write $P_j=[j/100,(j+1)/100]$. All moment intervals and probability
panels in the tables are closed; common endpoints are checked in both
adjacent pieces. The value $4097/10000$ is strictly below $\CE$.
The following are the exact low-moment input contracts. An own-score
ceiling bounds $C_n$, whereas a HIGH ceiling bounds only the named
high-frequency term in the compensated Fourier formula.
\begin{center}\small
\begin{tabular}{c|c|c|c|c}
actual $n$ & $\rho$ & initial radius & own-score ceiling & HIGH ceiling\\\hline
$15$--$19$ & $[1,11/10]$ & $23/10$ & $417/1000$ & $191/500$\\
$15$--$18$ & $[11/10,23/20]$ & $12/5$ & $827/2000$ & $409/1000$\\
$19$ & $[11/10,23/20]$ & $23/10$ & $827/2000$ & $409/1000$\\
$20$ & $[1,11/10]$ & $23/10$ & $831/2000$ & cellwise\\
$20$ & $[11/10,9/8]$ & $23/10$ & $103/250$ & cellwise\\
$21$--$24$ & $[1,11/10]$ & $23/10$ & $419/1000$ & cellwise\\
$21$--$24$ & $[11/10,9/8]$ & $12/5$ & $419/1000$ & cellwise\\
$25$--$30$ & $[1,11/10]$, $[11/10,6/5]$
 & actual support table & $434477/10^6$ & cellwise\\
$31$--$34$ & $[1,11/10]$, $[11/10,6/5]$
 & $47/20$, $13/5$ & $432313/10^6$ & cellwise
\end{tabular}
\end{center}
For an entry marked ``cellwise'', the compensated Fourier formula is
evaluated directly on each moment cell, without a separate scalar HIGH ceiling.
The floor-$25$ and floor-$31$ entries require the actual-index support
calculations in addition to the listed scalar ceilings. They do not
assert that support at index $25$ or $31$ transfers to another index.
The direct Gaussian criterion supplies each such calculation.

For $15\le n\le19$, support of radius $3$ is first established on
$1\le\rho\le33/25$. Supported Fourier bounds then establish the first
three own-score ceilings in the table; their hypotheses therefore do
not depend on the later, smaller support radii. The HIGH computation
uses this already acquired radius $3$. At index $15$ the retained
contact allowances in the two low bands are $219/2000$ and $57/1000$;
at indices $16$--$19$ they are zero. The exact Gaussian support tests
then establish the smaller radii used by the moment calculation.

At index $20$, the direct radius-$3$ result on $\rho\le6/5$ supplies
the support premise for the two supported Fourier own-score estimates
$831/2000$ and $103/250$. The two subsequent moment calculations
retain contact allowances $29/250$ and $23/500$, respectively.
They start without a numerical moment seed, perform eight refinements,
and export a ninth complete cover. There is no need to assume the
older coarse index-$20$ calculation in order to start this iteration.
The original coarse count contracts on $P_{37},\ldots,P_{40}$ do,
however, require a separate simultaneous-containment check against the
newly acquired constraints; their historical headers are not a proof
of that implication.

For $21\le n\le24$, the unrestricted auxiliary Fourier calculation
covers each actual index and gives $C_n<419/1000$. Its acceptance
checks this strict per-index ceiling, not merely the looser discovery
target $21/50$. Each band then uses eight refinements and a ninth
export, and all $100$ probability panels are retained. The direct
radius-$3$ theorem supplies the supported complement below $6/5$.

The high-moment alternatives are as follows:
\begin{center}\small
\begin{tabular}{c|c|c|c}
actual $n$ & moment interval & extra hypothesis & score ceiling\\\hline
$15$--$19$ & $[23/20,33/25]$ & $|X|\le3$ & $4097/10000$\\
$15$--$19$ & $[33/25,19/10]$ & none & $4097/10000$\\
$20$--$24$ & $[9/8,6/5]$ & $|X|\le3$ & $4097/10000$\\
$20$--$24$ & $[6/5,19/10]$ & none & $4097/10000$\\
$25$--$39$ & $[6/5,19/10]$ & none & $4097/10000$
\end{tabular}
\end{center}
The published decreasing Berry--Esseen majorant supplies $\rho\ge19/10$.
Every unrestricted row uses the full admissible auxiliary-moment
interval, not a selected value of that moment. Thus the last row
removes all need for the older index-$25$ band
$[6/5,13/10]$ and the subsequent high-moment resource chains. The
unused computations may be retained as supplementary evidence but
are not premises of this shorter proof.

\subsection{Finite acquisition and the common probability law}
The Fourier and contact bounds contain a few signed moments. Split
their complete admissible rectangle into closed cells. On a cell,
use outward bounds for the Fourier terms and for the signed Gaussian
maximum in the contact inequality. If the resulting upper bound is
negative, no maximizing witness lies in the cell. Otherwise retain
all necessary constraints established there. Taking minima of valid
upper bounds and maxima of valid lower bounds improves the same
witness's constraints. It does not require a contraction or convergence
of the iteration. A finite number of passes, followed by one complete
export pass, is sufficient.

The exported global defect upper bound is smaller than the positive
central-gap threshold. Hence zero is excluded from the support, and
sign conditioning is legitimate. In a probability panel $P_j$, the
resource calculation retains bounds for both central endpoints,
both outer endpoints, both signs of the phase, the signed moments,
and one common physical defect. Section~\ref{ext:count-filters} and
the resource-to-count calculation turn these into a simultaneous
count domain. In particular, if
\[
 D\le d,\qquad D\ge b+g d_{\rm in}+h\lambda,
\]
then $g d_{\rm in}+h\lambda\le d-b$. The triangle's original
variance intercept is retained even when a separate scalar estimate
further clips $\lambda$. Its two intercepts must not be independently
improved and then treated as a new joint triangle.

The floor-$25$ calculation is used at actual indices $26$--$30$,
and the floor-$31$ calculation at actual indices $32$--$34$.
Their invariant is $nD\le Nd$, where $N$ is the floor. This is stronger
than simply reusing $D\le d$. The count contract uses $D\le Nd/n$
and the corresponding noise coefficient $Nh/n$, with the actual
index in every binomial probability and Gaussian argument. The
actual-index support checks are separate inputs. These facts explain
why the reuse is valid and why it is limited to the stated finite
ranges.

The floor-$31$ low-moment calculation has two special local refinements
in the second band, on $P_{35}$ and $P_{36}$. The full global
acquisition supplies every other panel. At indices $32$--$34$ the
same weighted physical inequalities are used, with individually
checked directed count domains.

\subsection{The extra own-score bound at index 39}
Indices $35$--$39$ use actual-index INITIAL calculations, rather than
the floor-$31$ profile. Their input pairs (support radius, own-score
ceiling) are
\[
\begin{array}{c|ccccc}
 n&35&36&37&38&39\\\hline
 R&39/10&19/5&19/5&19/5&19/5\\
 c&53539/125000&1699/4000&84867/200000&211449/500000&41449/100000.
\end{array}
\]
Lemma~\ref{ext:first-failure25}, retaining its sharper expression
$C_n<41/100+4177/(13824n)$, implies the first four score ceilings.
The fifth requires an additional argument: the same-law local
comparison at actual index $39$ with
\[
 K=741/4000,\qquad q=33/100,\qquad C_{39}<2113/5000,
\]
excludes $39(C_{39}-\CE)\ge741/4000$. Its scalar entry conditions
follow from sharpness at $38$ and Lemma~\ref{ext:first-failure25}.
Since $\CE<40974/10^5$ and $741/(4000\cdot39)=19/4000$,
this gives $C_{39}<41449/10^5$, the required native ceiling.
Thus this positive-excess calculation is an internal prerequisite of
the selected sharp-$39$ calculation. Omitting it merely because
sharp-$39$ will later be available would be circular.

The actual-index traversals then cover every probability panel.
At $35$ they combine four local refinements with separate lower and
upper probability ranges; at $36$ only one interval needs a local
refinement. The remaining indices use their global INITIAL exports.
Appendix~\ref{ext:actual3539-acquisition} gives the complete union and
explains how the acquired bounds imply the $120$ surviving count
contracts, including their signed phase inequalities.

\subsection{Exact count coverage and assembly}
The following table specifies the surviving probability panels.
Every omitted panel must be explicitly excluded by its acquired
resource or scalar inequality. A failed resource evaluation is not
an exclusion. Multiple forests in one panel form a closed subdivision
in probability or another displayed coordinate; their union is
checked separately.
\begin{center}\small
\begin{tabular}{c|c|c|c}
$n$ & first low band & second low band & selected forests\\\hline
15 & $37$--$56$ & $33$--$44$ &32\\
16 & $37$--$56$ & $33$--$43$ &31\\
17 & $37$--$55$ & $33$--$43$ &30\\
18 & $37$--$54$ & $34$--$43$ &28\\
19 & $37$--$53$ & $34$--$42$ &26\\
20 & $37$--$58$ & $35$--$45$ &49\\
21 & $37$--$57$ & $35$--$44$ &31\\
22 & $37$--$56$ & $35$--$43$ &29\\
23 & $37$--$55$ & $35$--$43$ &28\\
24 & $37$--$54$ & $35$--$42$ &26\\
25 & $37$--$53$ & $32$--$42$ &52\\
26 & $37$--$53$ & Fourier &17\\
27 & $37$--$52$ & Fourier &16\\
28 & $37$--$52$ & Fourier &16\\
29 & $37$--$52$ & Fourier &16\\
30 & $37$--$51$ & Fourier &15\\
31--34 & $36$--$49$ & $30$--$40$ &25 per index\\
35 & \multicolumn{2}{c|}{$23$--$48$} &26\\
36 & \multicolumn{2}{c|}{$24$--$48$} &25\\
37 & \multicolumn{2}{c|}{$24$--$47$} &24\\
38 & \multicolumn{2}{c|}{$25$--$47$} &23\\
39 & \multicolumn{2}{c|}{$26$--$47$} &22
\end{tabular}
\end{center}
At $26\le n\le30$, the second-band resource survivors are
$P_{32},\ldots,P_{42}$; the selected supported Fourier calculation
checks $23$ closed moment strips, at each of the five actual indices,
using the acquired support and moment constraints on these panels.
\paragraph{Five panels excluded before the count calculation.}
The first-band count list shortens as $n$ increases because the upper
bound $D\le25d/n$ improves while the physical binary-asymmetry cost
does not decrease. The floor-$25$ calculation leaves the $17$ panels
$P_{37},\ldots,P_{53}$. Five of the $85$ resulting index--panel pairs
at $26\le n\le30$ are excluded as follows.

The complete acquisition gives these convenient outward bounds for the
same actual law:
\[
\begin{array}{c|c|c}
 j& d\ \text{in }D\le25d/n&\Theta\ \text{in }\theta\le\Theta\\\hline
 52&3/10&1/500\\
 53&5/16&1/2500.
\end{array}
\]
The physical variance bound is transferred by the argument in
Appendix~\ref{ext:finite-floor-transfer}; it is not a bound for an
unrelated optimizer at index $25$. Set $s=316227766/10^8<\sqrt{10}$
and $l=j/100$. Since $p\ge l>(4-s)/2\ge p_*$ and
$\sqrt{p(1-p)}\le1/2$, \eqref{ext:joint-physical-resource} gives
\[
 D\ge (12+4s)(1-\Theta)^2
            \left(l-\frac{4-s}{2}\right)^2
 >\begin{cases}251/1000,&j=52,\\38/125,&j=53.\end{cases}
\]
We used $(1-\theta)^{3/2}\ge(1-\Theta)^2$; the two final strict
comparisons are rational arithmetic. But on $P_{53}$ for
$n=27,28,29,30$ the upper bound gives
$D\le125/432<38/125$, and on $P_{52}$ for $n=30$ it gives
$D\le1/4<251/1000$. These contradictions exclude all five closed
panels, including their endpoints. The remaining $80$ selected cases
have full actual-index count certificates. Thus no missing certificate
is treated as an exclusion.

At index $20$, $P_{36}$ in the first band and $P_{34}$ in the second
are excluded by a strict Bernoulli-moment lower bound. These two
panels must not be silently omitted from the probability union.

For each selected forest, the source-to-domain join has four distinct
checks. First, the fresh resource bounds contain the same witness
on the entire original moment and probability domain. Second, all
scalar and phase constraints in the count contract are consequences
of those bounds, with upper bounds enlarged, lower bounds diminished,
and each signed phase quadratic dominated on its full phase interval.
Third, the native count calculation uses the actual $n$ and verifies
every closed leaf and every root count. Finally, the accepted leaves
and split panels form the complete closed domain. The common-count
identity then bounds the original convolution probability; it is not
an inference from the feasibility of a relaxation.

The direction of the phase comparison deserves emphasis. The count
kernel uses the maximum of its listed quadratic bounds, since an
exported list can represent alternative surviving moment cells.
If the acquired quadratics are $q_i(H)$ and the count quadratics are
$\widetilde q_j(H)$, the join verifies, for each $i$, some $j$ with
$q_i(H)\le\widetilde q_j(H)$ throughout the relevant signed interval.
It follows that
\[
 \max_i q_i(H)\le\max_j\widetilde q_j(H).
\]
The full-interval comparison checks the two endpoints and any interior
minimum of their quadratic difference. Replacing either maximum by a
minimum would require a different argument.

\begin{proposition}[Finite assembly criterion]\label{ext:finite1539-assembly}
Assume $C_{14}\le\CE$, the continuous support, Fourier, finite-moment,
resource, and count inequalities stated above, and the complete exact
arithmetic and source-to-domain joins for the displayed families.
Then $C_n\le\CE$ for every $15\le n\le39$.
\end{proposition}
\begin{proof}
If the conclusion failed, choose the least such $n$ and its positive
closed ratio maximizer. The moment tables cover its $\rho$. In a
Fourier region its score is already strictly below $\CE$, a
contradiction. In every remaining region the support and own-score
premises precede the moment acquisition; at $n=39$ include the
positive-excess prerequisite just described. The full exported
probability cover either excludes its panel or places the same
marginal in an accepted count domain. The count identity and the
strict leaf bound again give a score below $\CE$. This contradiction
proves the conclusion, and supplies the sharp predecessor for the
next actual index.
\end{proof}
